\begin{solution}{124}
    \begin{subsolution}
        当一束线偏振光沿 $x_{3}$ 轴进入晶体，且沿 $x_{1}$ 方向偏振时，进入晶体后可以分解为沿 $x_{1}'$ 和 $x_{2}'$ 方向的两个垂直偏振分量。据题给的折射率与电场强度的关系，相应的相位变化量为
        \begin{equation}
            \phi_{n_{x_{1}'}} = \frac{2 \pi}{\lambda_{0}} n_{x_{1}'} l = \frac{2 \pi l}{\lambda_{0}} (n_{\mathrm{o}} - \frac{1}{2} n_{\mathrm{o}}^{3} \gamma E)
            \label{eq:1.s124}
        \end{equation}
        \begin{equation}
            \phi_{n_{x_{2}'}} = \frac{2 \pi}{\lambda_{0}} n_{x_{2}'} l = \frac{2 \pi l}{\lambda_{0}} (n_{\mathrm{o}} + \frac{1}{2} n_{\mathrm{o}}^{3} \gamma E)
            \label{eq:2.s124}
        \end{equation}
        这两束光穿过晶体后的相位差为
        \begin{equation}
            \Delta \phi = \phi_{n_{x_{2}'}} - \phi_{n_{x_{1}'}} = \frac{2 \pi}{\lambda_{0}} l n_{\mathrm{o}}^{3} \gamma E = \frac{2 \pi}{\lambda_{0}} n_{\mathrm{o}}^{3} \gamma V
            \label{eq:3.s124}
        \end{equation}
        式中 $V = E l$。入射光强为
        \begin{equation}
            I_{\mathrm{i}} = k |\mathbf{E}|^{2} = k (|\mathbf{E}_{x_{1}'}|^{2} + |\mathbf{E}_{x_{2}'}|^{2}) = 2 k E_{0}^{2}
            \label{eq:4.s124}
        \end{equation}
        式中 $k$ 是与电场强度无关的常量。当光通过长度为 $l$ 的晶体后，在输出面的 $x_{1}'$ 与 $x_{2}'$ 方向上的电场为
        \begin{equation}
            \left\{ \begin{array}{l} E_{x_{1}'}(l) = E_{0} \cos (\omega t + \phi_{n x_{1}'}) \\ E_{x_{2}'}(l) = E_{0} \cos (\omega t + \phi_{n x_{2}'}) \end{array} \right.
            \label{eq:5.s124}
        \end{equation}
        出射光强为
        \begin{equation}
            I_{\mathrm{f}} = k |\mathbf{E}(l)|^{2} = k (|\mathbf{E}_{x_{1}'}(l)|^{2} + |\mathbf{E}_{x_{2}'}(l)|^{2}) = 2 k E_{0}^{2} = I_{\mathrm{i}}
            \label{eq:6.s124}
        \end{equation}
        可见，当光通过长度为 $l$ 的晶体后，光强未变。

        由半波电压 $V_{\pi}$ 的定义有
        \begin{equation}
            \Delta \phi = \frac{2 \pi}{\lambda_{0}} n_{\mathrm{o}}^{3} \gamma V_{\pi} = \pi
            \label{eq:7.s124}
        \end{equation}
        由 \eqref{eq:7.s124} 式得
        \begin{equation}
            V_{\pi} = \frac{\lambda_{0}}{2 n_{\mathrm{o}}^{3} \gamma}
            \label{eq:8.s124}
        \end{equation}
    \end{subsolution}
    \begin{subsolution}
        不经过 1/4 波片，从 $\mathrm{P}_{2}$ 射出的光场强度为 $x_{1}'$ 与 $x_{2}'$ 的 $x_{2}$ 分量之和
        \begin{align}
            E_{x 2} & = -E_{x_{1}'}(l) \cos 45^{\circ} + E_{x_{2}'}(l) \cos 45^{\circ} \nonumber \\
            & = \frac{E_{0}}{\sqrt{2}} [-\cos (\omega_{0} t + \phi_{n_{x_{1}'}}) + \cos (\omega_{0} t + \phi_{n_{x_{2}'}})] \nonumber \\
            & = -\frac{2 E_{0}}{\sqrt{2}} \sin (\frac{\Delta \phi}{2}) \sin (\omega_{0} t + \frac{\phi_{n_{x_{2}'}} + \phi_{n_{x_{1}'}}}{2})
            \label{eq:9.s124}
        \end{align}
        输出光强为
        \begin{equation}
            I_{\mathrm{f}} = k \left(\frac{2 E_{0}}{\sqrt{2}} \sin \frac{\Delta \phi}{2}\right)^{2} = 2 k E_{0}^{2} \sin^{2} \frac{\Delta \phi}{2}
            \label{eq:10.s124}
        \end{equation}
        由 \eqref{eq:3.s124} \eqref{eq:4.s124} \eqref{eq:8.s124} \eqref{eq:10.s124} 式得，当 $\frac{\pi V_{\mathrm{m}}}{2 V_{\pi}} << 1$，光强的透过率为
        \begin{equation}
            T \equiv \frac{I_{\mathrm{f}}}{I_{\mathrm{i}}} = \sin^{2} \left(\frac{\Delta \phi}{2}\right) = \sin^{2} \left(\frac{\pi V}{2 V_{\pi}}\right) = \sin^{2} \frac{\pi (V_{\mathrm{m}} \sin \omega_{\mathrm{m}} t)}{2 V_{\pi}} \approx \left[ \frac{\pi (V_{\mathrm{m}} \sin \omega_{\mathrm{m}} t)}{2 V_{\pi}} \right]^{2}
            \label{eq:11.s124}
        \end{equation}
        $T$ 与 $V^{2} = (V_{\mathrm{m}} \sin \omega_{\mathrm{m}} t)^{2}$ 成正比。

        加上 1/4 波片，$x_{1}'$ 与 $x_{2}'$ 传输的光波分量之间有一个固定的 $\pi/2$ 相位差，由 \eqref{eq:11.s124} 式有
        \begin{align}
            T & = \sin^{2} \frac{\Delta \varphi + \frac{\pi}{2}}{2} = \sin^{2} \left[ \frac{\pi}{4} + \frac{\pi}{2 V_{\pi}} (V_{\mathrm{m}} \sin \omega_{\mathrm{m}} t) \right] \nonumber \\
            & = \frac{1}{2} \left\{1 - \cos [2 (\frac{\pi}{4} + \frac{\pi V_{\mathrm{m}} \sin \omega_{\mathrm{m}} t}{2 V_{\pi}})] \right\} = \frac{1}{2} (1 + \sin \frac{\pi V_{\mathrm{m}} \sin \omega_{\mathrm{m}} t}{V_{\pi}}) \nonumber \\
            & \approx \frac{1}{2} (1 + \frac{\pi}{V_{\pi}} V_{\mathrm{m}} \sin \omega_{\mathrm{m}} t)
            \label{eq:12.s124}
        \end{align}
        $T$ 与 $V = V_{\mathrm{m}} \sin \omega_{\mathrm{m}} t$ 成线性关系。

        可见，未加 1/4 波片时，光强的透过率与调制电压信号 $V = V_{\mathrm{m}} \sin \omega_{\mathrm{m}} t$ 呈非线性关系；添加 1/4 波片的作用是保证系统的工作电压处于线性工作区域，使得光强的透过率和调制电压信号 $V = V_{\mathrm{m}} \sin \omega_{\mathrm{m}} t$ 呈现线性关系。
    \end{subsolution}
    \begin{subsolution}
        实际加载至电光晶体上的电压 $V_{\mathrm{c}}$ 为
        \begin{equation}
            V_{\mathrm{c}} = \frac{\left(\frac{1}{1/R + i \omega_{\mathrm{m}} C}\right)}{R_{\mathrm{m}} + \frac{1}{1/R + i \omega_{\mathrm{m}} C}} V
            \label{eq:13.s124}
        \end{equation}
        由 \eqref{eq:13.s124} 式得，$V_{\mathrm{c}}$ 的模为
        \begin{align}
            |V_{\mathrm{c}}| & = \frac{R}{\sqrt{(R_{\mathrm{m}} + R)^{2} + R_{\mathrm{m}}^{2} R^{2} C^{2} \omega_{\mathrm{m}}^{2}}} |V| \nonumber \\
            & \approx \frac{1}{\sqrt{1 + C^{2} R_{\mathrm{m}}^{2} \omega_{\mathrm{m}}^{2}}} |V|
            \label{eq:14.s124}
        \end{align}
        当 $\omega_{\mathrm{m}} >> (R_{\mathrm{m}} C)^{-1}$ 时，
        \begin{equation}
            |V_{\mathrm{c}}| << |V|
        \end{equation}
        电压基本上落在 $R_{\mathrm{m}}$ 上，电光晶体调制效率极低。
    \end{subsolution}
    \begin{subsolution}
        并联电感 $L$ 和分流电阻 $R_{\mathrm{L}}$（即将图中的开关都闭合）后，调制器系统的有效阻抗满足
        \begin{equation}
            \frac{1}{Z} = \frac{1}{R_{\mathrm{L}}} + \frac{1}{i \omega_{\mathrm{m}} L} + \frac{1}{-i / (\omega_{\mathrm{m}} C)} = \frac{1}{R_{\mathrm{L}}} - i \left(\frac{1}{\omega_{\mathrm{m}} L} - \omega_{\mathrm{m}} C\right)
            \label{eq:15.s124}
        \end{equation}
        由 \eqref{eq:15.s124} 式得
        \begin{equation}
            |Z| = \frac{1}{\sqrt{\frac{1}{R_{\mathrm{L}}^{2}} + \left(\omega_{\mathrm{m}} C - \frac{1}{\omega_{\mathrm{m}} L}\right)^{2}}}
            \label{eq:16.s124}
        \end{equation}
        由 \eqref{eq:16.s124} 式可知，当 $\omega_{\mathrm{m}} = 1/\sqrt{LC}$，调制器阻抗为最大值
        \begin{equation}
            Z_{\max} = R_{\mathrm{L}}
            \label{eq:17.s124}
        \end{equation}
        因此，可通过调节电感来使得调制器阻抗达到最大值。特别的，当 $R_{\mathrm{L}} >> R_{\mathrm{m}}$ 时有
        \begin{equation}
            |V_{\mathrm{c}}| \approx |V|
            \label{eq:18.s124}
        \end{equation}
        调制电压绝大部分加在晶体上，从而保证了调制效率。
    \end{subsolution}
\end{solution}